\documentclass[11pt,a4paper]{article}
\pdfinfoomitdate=1\pdftrailerid{}
\usepackage[margin=25mm]{geometry}\usepackage[T1]{fontenc}\usepackage{mathpazo,hyperref,xurl}
\hypersetup{hidelinks}\setlength{\parindent}{0pt}\setlength{\parskip}{7pt}\setlength{\emergencystretch}{2em}
\begin{document}
4 November 2026

Dear SoftwareX Editors,

Please consider \emph{WorkCert: Reconstruction error metadata for force--displacement tables} as an Original Software Publication describing WorkCert 0.3.0.

WorkCert combines reduced force--displacement tables with continuous work-error metadata, compensated integration, local window queries and checked persistence. A public-data consumer workflow requests finer tables from two source records to identify the leader at reported masses, while a mass-rounding scenario shows when better mass information is needed. The same interface inverts a prescribed work capacity into a stroke interval. A complete cyclic-bending history illustrates directed branch composition and explicit detection of an unmet point-capped budget.

Evaluation includes constant and affine scale families, strong-cancellation windows, corrupted metadata and overflow regressions; geometric selectors and objective-matched dynamic programming contextualize selection. The accompanying material contains source data, preparation and plotting programs, numerical outputs, installation instructions and tests. The fixed software source is \url{https://github.com/changs1723/WorkCert/tree/v0.3.0}. The software archive is \url{https://doi.org/10.5281/zenodo.23934502}.  The final article sources, figures and build inputs are fixed at \url{https://github.com/changs1723/WorkCert/tree/article-v0.3.0-final}. 

\textbf{Declaration of generative AI and AI-assisted technologies in the manuscript preparation process.} OpenAI Codex (OpenAI; local CLI version 0.166.0-alpha.5.1) assisted with research planning, literature and public-data collection, implementation and debugging, numerical analysis and verification, figure preparation, drafting, language editing and typesetting. The authors are responsible for the research design, interpretation, software, results and final manuscript.

The authors are Saisai Hou, Yunzhi Zhu and Sen Zhang, all affiliated with the Department of Public Basic Courses, Nanjing University of Industry Technology. Sen Zhang is the corresponding author. The work is supported by Grant No. 25KJD110002 of the General Program of Natural Science Research (Basic Science) in Higher Education Institutions of Jiangsu Province.

Sincerely,

Sen Zhang\\Corresponding author\\\href{mailto:szhang@niit.edu.cn}{szhang@niit.edu.cn}
\end{document}
